\MakeAbstract{
    Unix V6++ 是面向操作系统课程实践的教学内核，重构前系统运行于 IA-32 传统 PC 平台，内核主体采用 C++ 实现。将其迁移到 RISC-V 后，启动入口、地址空间、中断模型、设备访问路径和用户程序格式均需重新适配。本文研究 Unix V6++ 的 RISC-V 移植与 Rust 重构，在保留 Unix V6 风格系统调用、进程语义、文件接口和 Shell 使用方式的前提下，重写平台相关层，使系统能够在 QEMU virt 环境中启动并运行用户态程序。

    当前系统面向 QEMU virt 平台和 OpenSBI 固件环境，内核使用 Rust 编写，启动、Trap 入口和上下文切换保留少量汇编。论文重点讨论 RISC-V Trap 与系统调用框架、进程控制块、调度、fork/exec/wait/exit、ELF64/RISC-V 用户程序加载、信号处理以及 TTY/Shell 交互验证；文件系统和内存管理主要作为系统运行支撑条件处理。相较重构前的 main 分支，当前系统把 Unix V6++ 从 IA-32/C++/PC 设备模型迁移到 RISC-V/Rust/QEMU virt 平台，并形成了可构建、可运行、可交互测试的教学内核。Rust 的所有权、生命周期、Drop 和 RAII 风格封装用于表达页资源、内核栈、文件引用、inode 引用和中断保护等对象的释放边界，使资源归属和回收位置更易追踪。
}{Unix V6++；RISC-V；Rust；操作系统移植；教学操作系统}

\MakeAbstractEng{
    Unix V6++ is an educational operating system kernel used in operating system courses. The pre-refactoring system targets the IA-32 PC platform and is mainly implemented in C++. Porting it to RISC-V requires new boot code, address-space support, trap handling, device access, and user program ABI support. This thesis studies the RISC-V porting and Rust-based refactoring of Unix V6++, with Unix V6-style system calls, process semantics, file interfaces, and shell behavior preserved as the main teaching interface.

    The refactored system runs on the QEMU virt machine with OpenSBI. The kernel is written in Rust, with a small amount of assembly for boot code, trap entry, and context switching. The thesis focuses on the RISC-V trap and system call framework, process control block, scheduling, fork/exec/wait/exit, ELF64/RISC-V program loading, signal handling, and TTY/shell validation. File system and memory management are treated as supporting parts of the runtime environment. Compared with the pre-refactoring main branch, the current system migrates Unix V6++ from an IA-32/C++/PC device model to a RISC-V/Rust/QEMU virt platform and provides a buildable, runnable, and interactively testable educational kernel. Rust ownership, lifetimes, Drop, and RAII-style wrappers describe the release boundaries of page resources, kernel stacks, file references, inode references, and interrupt guards, making resource ownership and cleanup paths easier to follow.
}{Unix V6++; RISC-V; Rust; operating system porting; educational operating system}
